\section*{零基础读法与符号准备}
本册逐页保留原课件图像，再解释其数学含义。原图出现笔误时，在正文中明确纠正；截图与讲解的不同不代表遗漏。公式阅读先说中文：$\forall$ 是“对每一个”，$\exists$ 是“至少存在一个”，$\Rightarrow$ 是“能够推出”，$\Leftrightarrow$ 要证明两个方向。

\paragraph{最小工具箱}$x\in S$ 表示属于集合；$A\subseteq B$ 表示 $A$ 中每个元素都在 $B$ 中；$\bigcap_n A_n$ 要同时属于所有集合，$\bigcup_n A_n$ 只须属于至少一个。$\sup$ 是最小上界，$\inf$ 是最大下界，二者不保证属于原集合。绝对值 $|x-y|$ 表示一维距离。

\paragraph{怎样写证明}先抄清假设与目标，再把目标翻译成定义。证明“对所有”时先任取对象；证明“存在”时要构造对象并检查条件；反证法是暂时假定目标不成立并导出矛盾。否定“对每个 $\eps$ 存在 $K$”时，量词逐个反转，不能只在最后加一个“不”。

\paragraph{学习安排}每次阅读一个证明后，合上正文，回答：候选对象从哪里来？每个假设在哪里用到？删去一个假设有何反例？能回答这三问，再练习完整书写。原图中的英文术语可辅助对照，不要求先背英文定义。


\Slide{1}{本讲要解决什么问题}
这份补充课件从“实数是什么”讲到有限维空间中的完备性与紧性。不要把几十条定理当作互不相干的结论：主线是实数轴没有缺口，因此逼近过程有落脚点，进而能从无限信息中提取有限信息。
\paragraph{阅读目标}第一遍会读定义，第二遍能解释每个假设，第三遍闭卷写证明。最重要的链条是：确界存在 $\Rightarrow$ 单调有界收敛 $\Rightarrow$ 闭区间套 $\Rightarrow$ 有界数列有收敛子列 $\Rightarrow$ 柯西准则。有限覆盖是同一完备性的另一个面貌。
本页是标题页，没有需要补证的数学命题。后面的“原页”编号均指这份 PDF 的物理页码。

\Slide{2}{从自然数逐步构造实数}
$\N=\{1,2,\ldots\}$ 用来计数；加入 $0$ 与负数得到 $\Z$；允许非零整数作分母得到 $\Q$。但 $\Q$ 中没有平方等于 $2$ 的数，说明只会四则运算还不足以容纳所有极限。
\paragraph{构造的直觉}用 $1,1.4,1.41,1.414,\ldots$ 逐渐逼近 $\sqrt2$。每一项都是有理数，目标却不在 $\Q$ 中。柯西序列的方法把这样的“越来越稳定的逼近”作为新数。
\paragraph{准确补充}严格地说，实数不是所有有理柯西序列本身，而是这些序列的\emph{等价类}：$(a_n)\sim(b_n)$ 当 $a_n-b_n\to0$。不同逼近过程可能代表同一个数。Dedekind 分割则用目标左侧和右侧的有理数来定位目标；本课使用性质而不要求从头建设整个数系。

\Slide{3}{实数是完备有序域}
“域”保证加减乘除（除数非零）合法；“有序”保证能比较大小且运算保持相容；“完备”保证有界的逼近不会落进一个缺口。三个层次缺一不可。
\paragraph{为什么必须强调完备}有理数也是有序域，却不是完备有序域。因而证明中一旦说“取 $\sup S$”或“这个柯西序列有极限”，实际上就用了比四则运算更强的性质。
\paragraph{证明意识}公理是起点，不要求在同一体系内再证明。选择确界公理后，单调收敛、闭区间套等才是要从它推出的定理；不能互相引用绕成一圈。

\Slide{4}{加法群：能加、能撤销}
封闭性要求 $a+b$ 仍在集合里；结合律只改变括号；单位元 $0$ 加上去不改变数；逆元 $-a$ 撤销加上 $a$ 的效果。交换律 $a+b=b+a$ 是另外一个条件，满足它才叫交换群。
\paragraph{两个小证明}若 $e,e'$ 都是单位元，则 $e=e+e'=e'$，故单位元唯一。若 $b,c$ 都是 $a$ 的加法逆元，则
\[b=b+0=b+(a+c)=(b+a)+c=0+c=c.\]
每个等号依次使用单位元、逆元、结合律、逆元、单位元。这就是“由公理证明”的写法。
\paragraph{例子与反例}$\Z$ 对加法是群；$\N$ 对加法不是群，因为缺 $0$ 和负数。判断群时必须同时说明集合与运算。

\Slide{5}{乘法群：为什么排除零}
乘法单位元是 $1$，乘法逆元是 $a^{-1}=1/a$。$\R\setminus\{0\}$ 对普通乘法是交换群，$\R$ 对普通乘法不是群，因为 $0b=0$ 不可能等于 $1$。
\paragraph{逆元唯一}若 $ab=ba=1$ 且 $ac=ca=1$，则 $b=b(ac)=(ba)c=c$。这里不能把“逆元”与“负数”混淆：$2$ 的加法逆元是 $-2$，乘法逆元是 $1/2$。
\paragraph{自检}$\Z\setminus\{0\}$ 对乘法不是群，因为 $2$ 没有整数逆元；正实数对乘法是群，因为正数的乘积及倒数仍为正数。

\Slide{6}{一般群与消去律}
本页把加法、乘法统一写成运算 $\cdot$。符号不意味着一定是普通乘法；例如置换的复合也是群运算。$e$ 是单位元，$a^{-1}$ 是相对于该运算的逆元。
\paragraph{左消去律的证明}若 $a\cdot b=a\cdot c$，两边左乘 $a^{-1}$，由结合律得
\[(a^{-1}\cdot a)\cdot b=(a^{-1}\cdot a)\cdot c,\qquad e\cdot b=e\cdot c,\]
因此 $b=c$。右消去同理，但应从右侧乘逆元。在非交换群中不能任意调换因子。
\paragraph{层次}封闭性、结合律、单位元、逆元是定义；唯一性和消去律是由定义推出的性质。考试写证明时要分清两者。

\Slide{7}{环：把加法和乘法放在一起}
本课采用有乘法单位元的环定义：加法构成交换群；乘法封闭、有结合律和单位元；两种运算满足左右分配律。一般环的乘法不必交换，也不要求非零元素都有倒数。
\paragraph{基本推论}由 $a(0+0)=a0+a0$ 及 $0+0=0$，得 $a0=a0+a0$。两边加上 $-(a0)$ 得 $a0=0$。同理 $0a=0$。
\paragraph{例子}$\Z$ 是环，但 $2$ 在 $\Z$ 中不可逆；实方阵集合也是环，矩阵乘法一般不交换。不要把这里表示任意环的字母 $R$ 自动读成实数集 $\R$。

\Slide{8}{交换环比一般环多了什么}
本页比上一页增加乘法交换律 $ab=ba$，而加法交换律在环定义中已经有了。整数、有理数、实数、多项式都提供熟悉的例子。
\paragraph{不要误用约分}在一般交换环里，$ab=ac$ 且 $a\ne0$ 未必推出 $b=c$。模 $6$ 的余数环中 $2\cdot1=2\cdot4$，但 $1\ne4$。失败来自 $2\cdot3=0$ 这样的零因子。
\paragraph{考试检查}写“除以 $a$”前，应确认 $a$ 有乘法逆元；只有非零还不一定足够。下一页的域正是让每个非零元素都可逆的结构。

\Slide{9}{域：合法四则运算的基础}
域是满足 $0\ne1$ 的交换环，并且每个非零元素有乘法逆元。因此 $a/b$ 的含义是 $a\cdot b^{-1}$，条件 $b\ne0$ 不能省略。
\paragraph{无零因子的证明}设 $ab=0$。若 $a=0$，结论成立；若 $a\ne0$，左乘 $a^{-1}$ 得 $b=a^{-1}0=0$。因此域中 $ab=0$ 必有一个因子为零。这一性质使许多代数约分合理。
\paragraph{辨认}$\Q,\R,\mathbb C$ 都是域，$\Z$ 不是域。本页把三组条件称作 ring axioms 属于不够精确的表述；含逆元条件后已是域公理。末行变量所属集合应为 $F$。

\Slide{10}{整环、除环和域}
整环是无零因子的含幺交换环；除环要求每个非零元素可逆，但乘法不必交换；交换的除环就是域。因而“整环”不等于“域”，$\Z$ 是直接反例。
\paragraph{关系要这样读}$\Z\subseteq\Q\subseteq\R\subseteq\mathbb C$ 是集合包含，并且普通四则运算在这些集合间相容。所有列出的数系都是环，只有后面三个是域。
\paragraph{小练习及答案}证明整环中非零元素可消去。若 $ab=ac$，则 $a(b-c)=0$；由 $a\ne0$ 和无零因子得 $b-c=0$。注意这里证明了消去律，却没有创造出一个不存在的整数 $1/2$。

\Slide{11}{向量空间中的两类对象}
向量空间包含向量 $v\in V$，其标量来自域 $F$。向量相加是 $V\times V\to V$；数乘是 $F\times V\to V$。例如 $V=\R^2,F=\R$，$2(1,3)=(2,6)$。
\paragraph{最易混淆处}数乘的第一个输入是数，第二个是向量；通常没有定义两个向量相乘仍为向量。经济学中的商品束 $(x_1,\ldots,x_N)$ 就可放在 $\R^N$ 中。
\paragraph{例子}所有连续实函数也构成向量空间，运算逐点定义：$(f+g)(t)=f(t)+g(t)$。因此“向量”不一定是画在平面上的箭头，关键是满足下页的运算法则。

\Slide{12}{向量空间八条公理如何使用}
前四条说明向量加法是交换群；后四条说明数乘与域运算配合。所有公理都要注意变量类型：$v,v_1,v_2\in V$，$a,a_1,a_2\in F$。
\paragraph{推导 $0v=0_V$}利用分配律，$0v=(0+0)v=0v+0v$，在向量加法群中消去一个 $0v$ 即得。类似地，$a0_V=a(0_V+0_V)=a0_V+a0_V$，所以 $a0_V=0_V$。又因为 $v+(-1)v=(1-1)v=0_V$，故 $(-1)v=-v$。
\paragraph{原页勘误}第(1)条把向量写成属于 $F$，应为属于 $V$。两个“零”分别属于域和向量空间，证明时写 $0_V$ 可避免混淆。

\Slide{13}{用正数集合定义大小}
$P$ 是严格正数集。三歧性说一个数恰好为零、为正或其相反数为正；正数之和、积仍为正。定义 $a<b$ 等价于 $b-a\in P$，把大小关系转成正数运算。
\paragraph{传递性的证明}若 $a<b$ 且 $b<c$，则 $b-a,c-b\in P$。正数相加仍正，故 $(b-a)+(c-b)=c-a\in P$，即 $a<c$。
\paragraph{不等式法则}由 $(b+z)-(a+z)=b-a$，同时加 $z$ 不改变大小。若 $c>0$，则 $bc-ac=c(b-a)>0$；若 $c<0$，乘上 $c$ 会反向。这里 $\R_+=[0,\infty)$，$\R_{++}=(0,\infty)$，须服从老师约定。

\Slide{14}{有序域与复数的区别}
有序域是域结构和满足相容条件的全序合在一起。$\Q$ 与 $\R$ 都是有序域，但只有 $\R$ 完备。
\paragraph{为什么复数不能这样排序}在有序域中，任意非零 $a$ 的平方为正：$a>0$ 时直接用正数乘法；$a<0$ 时 $a^2=(-a)^2>0$。因此 $1=1^2>0$，而 $i^2=-1$ 不可能也为正。故不能给 $\mathbb C$ 配上与域运算相容的实数式大小关系。
\paragraph{符号说明}课件的 $0_\oplus$ 是加法单位元，$0_\otimes$ 表示乘法单位元，后者通常写作 $1$；这不表示实数的 $0$ 与 $1$ 相等。

\Slide{15}{Dedekind 分割：数轴没有断口}
把全部实数分为非空的左部 $A$ 和右部 $B$，每个左部数都小于每个右部数。完备性说存在唯一分界点 $c$：小于它的都在 $A$，大于它的都在 $B$，而 $c$ 自己必须分到一边。
\paragraph{唯一性证明}若 $c<d$ 都能作分界点，取中点 $m=(c+d)/2$。由 $m>c$ 得 $m\in B$，由 $m<d$ 得 $m\in A$，矛盾。因为原条件使 $A\cap B=\varnothing$。
\paragraph{为什么两部分非空}若 $A=\R,B=\varnothing$，就没有有限分界点。若在 $\Q$ 中用平方小于或大于 $2$ 的条件切割，缺失的 $\sqrt2$ 正说明相同结论不适用于 $\Q$。

\Slide{16}{分界点究竟属于哪边}
本页把分割公理写成“左边有最大值，或右边有最小值”。若 $c\in A$，则 $c=\max A$、$A=(-\infty,c]$；若 $c\in B$，则 $c=\min B$、$B=[c,\infty)$。
\paragraph{为什么不能两边同时有端点}若 $a=\max A$ 且 $b=\min B$，分割条件给 $a<b$，但中点 $(a+b)/2$ 既不能在 $A$ 也不能在 $B$，与并集是整个实数轴矛盾。
$A^c$ 表示相对于 $\R$ 的补集，即 $\R\setminus A$。这页没有引入另一种完备性，只是上一页的等价说法。

\Slide{17}{上界、下界、上确界、下确界}
对集合 $S$，上界 $u$ 要满足“每一个 $x\in S$ 都有 $x\le u$”；上确界 $b$ 除了是上界，还必须小于等于任何其他上界。下界与下确界把不等号反过来。
\paragraph{例子}$S=(0,1)$ 的上界是 $[1,\infty)$，最小的上界是 $1$，故 $\sup S=1$；其下界是 $(-\infty,0]$，故 $\inf S=0$。但 $S$ 没有最大、最小元素，因为 $0,1$ 不在其中。
\paragraph{唯一性}若 $b,c$ 都是上确界，由 $c$ 是上界得 $b\le c$，交换角色得 $c\le b$，所以 $b=c$。证明一个候选数是确界时，必须检查“界”与“最优的界”两部分。

\Slide{18}{确界存在定理的假设}
非空且有上界的实数集合有有限实数上确界；非空且有下界的实数集合有有限实数下确界。“非空”“有界的方向”“实数”都是实质条件。
\paragraph{由上确界推出下确界}令 $-S=\{-x:x\in S\}$。若 $S$ 非空且有下界，则 $-S$ 非空且有上界，设 $b=\sup(-S)$。对每个 $x\in S$，$-x\le b$，故 $x\ge-b$；若 $d$ 是 $S$ 的下界，则 $-d$ 是 $-S$ 的上界，故 $b\le-d$，即 $d\le-b$。所以 $\inf S=-\sup(-S)$。
\paragraph{考试提醒}确界不一定在集合中。“存在确界”不能改写成“存在最大值”。

\Slide{19}{由分割公理证明上确界存在}
\paragraph{构造}令 $B$ 为 $S$ 的全部上界，$A=\R\setminus B$。有上界使 $B\ne\varnothing$；选 $s\in S$，$s-1$ 不是上界，故 $A\ne\varnothing$。
\paragraph{验证分割}若 $a\in A,b\in B$ 而 $a\ge b$，则所有 $s\in S$ 满足 $s\le b\le a$，使 $a$ 也成上界，矛盾。因此 $a<b$，分割公理给分界点 $c$。
\paragraph{证明端点属于上界侧}若 $c\notin B$，有 $s\in S$ 满足 $s>c$。中点 $m=(c+s)/2$ 大于 $c$，应属于 $B$；但 $m<s$，又不是上界，矛盾。故 $c\in B$。所有 $b\in B$ 都有 $c\le b$，所以 $c=\min B=\sup S$。
关键技巧：用“全部上界”形成新集合，不是直接把原集合随意切开。

\Slide{20}{由上确界存在证明分割公理}
设 $A,B$ 是一个分割。$A$ 非空且任取 $b\in B$ 都是其上界，故可取 $c=\sup A$。
\paragraph{验证分界性质}若 $x<c$ 而 $x\in B$，则 $x$ 是 $A$ 的上界，违背 $c$ 是最小上界，所以 $x\in A$。若 $x>c$，由于 $c$ 是 $A$ 的上界，$x$ 不可能在 $A$，只能在 $B$。
\paragraph{唯一性补全}两个不同分界点之间的中点会被同时分入 $A$ 和 $B$，矛盾。故确界公理能推出完整的分割公理。考试中“等价”必须写出两个方向，本页与前页缺一不可。

\Slide{21}{由分割公理证明下确界存在}
令 $A$ 为 $S$ 的全部下界，$B=\R\setminus A$。由下有界，$A\ne\varnothing$；取 $s\in S$，$s+1$ 不是下界，故 $B\ne\varnothing$。若 $a\in A,b\in B$ 且 $a\ge b$，则 $b\le a\le s$ 对所有 $s\in S$ 成立，矛盾，故 $a<b$。
设分界点为 $c$。如果 $c\in B$，则有 $s\in S$ 满足 $s<c$。中点 $m=(s+c)/2$ 小于 $c$，应在 $A$，但 $m>s$ 又不可能是下界。因此 $c\in A$，且所有下界 $a\le c$，故 $c=\inf S$。
\paragraph{记忆方法}上确界证明中选择“全部上界”的最小者；这里选择“全部下界”的最大者。中点总是用来挤出与分界方向冲突的数。

\Slide{22}{由下确界存在恢复分割}
$B$ 非空且任一 $a\in A$ 是其下界，取 $c=\inf B$。若 $x<c$，$x$ 不可能属于 $B$，否则违反 $c$ 是下界，所以 $x\in A$。若 $x>c$ 而 $x\in A$，则 $x$ 是 $B$ 的一个\emph{大于} $c$ 的下界，违背 $c$ 为最大下界，因此 $x\in B$。
\paragraph{原页勘误}末段英文写成 ``less than the greatest lower bound''，应为 ``greater than''。证明中的矛盾恰好来自找到一个更大的下界。
\paragraph{练习}把每个“不能”改写成反证的完整一句。这样能检查究竟使用了“$c$ 是下界”还是“$c$ 是最大下界”。

\Slide{23}{确界的逼近刻画}
$b=\sup S$ 当且仅当：所有 $x\in S$ 满足 $x\le b$，而且每给 $\eps>0$，都能找到 $x\in S$ 使 $b-\eps<x\le b$。下确界类似：$\ell\le x<\ell+\eps$。
\paragraph{直觉}集合可能碰不到边界，但一定能从内部任意接近边界。$\eps$ 是任意精度要求；找到的 $x$ 可以随 $\eps$ 改变，不要求同一个 $x$ 满足所有精度。
\paragraph{量词陷阱}“每个 $\eps$ 都有一个 $x$”写作 $\forall\eps>0\ \exists x\in S$；若颠倒成 $\exists x\ \forall\eps>0$，便强迫 $x=b$，这就成了最大值，强得多。后两页分别证明这两个等价刻画。

\Slide{24}{上确界逼近刻画的双向证明}
\paragraph{必要性}设 $b=\sup S$。它当然是上界。如果某个 $\eps>0$ 下找不到 $x>b-\eps$，那么所有 $x\in S$ 都满足 $x\le b-\eps$，说明 $b-\eps$ 是更小上界，矛盾。
\paragraph{充分性}设 $b$ 是上界且有逼近性质。任取上界 $u$，要证 $b\le u$。若 $u<b$，令 $\eps=b-u>0$。逼近性质给 $x\in S$ 满足 $x>b-\eps=u$，违反 $u$ 是上界。因此 $b\le u$。
\paragraph{如何想到选精度}想让 $b-\eps$ 恰好落在反证假设中的 $u$，所以解方程得到 $\eps=b-u$。证明中不是“随便取一个小数”，而是围绕目标构造。

\Slide{25}{下确界逼近刻画的双向证明}
\paragraph{必要性}若 $\ell=\inf S$，它是下界。若某个 $\eps>0$ 使所有 $x\in S$ 都满足 $x\ge\ell+\eps$，那么 $\ell+\eps$ 是更大的下界，矛盾。
\paragraph{充分性}若 $\ell$ 是下界并能从上方逼近，任取下界 $d$。假设 $d>\ell$，令 $\eps=d-\ell$，就有 $x\in S$ 满足 $x<\ell+\eps=d$，违背下界性质。故每个下界 $d\le\ell$。
\paragraph{一眼识别使用时机}需要构造 $x_n\in S$ 逼近下确界时，逐次取 $\eps=1/n$，得到 $\ell\le x_n<\ell+1/n$，于是 $x_n\to\ell$。存在逼近数列不意味着下确界属于 $S$。

\Slide{26}{单调与有界是两种不同限制}
数列是按次序排列的无限列表 $(x_1,x_2,\ldots)$。递增在本课指 $x_n\le x_{n+1}$，允许相等；严格递增才要求 $<$。有界指存在一个固定 $M$，使每个 $n$ 都有 $|x_n|\le M$。
\paragraph{例子}$1-1/n$ 递增且上有界；$1/n$ 递减且下有界；$n$ 递增但上无界；$(-1)^n$ 有界但不单调。后两例说明只满足一半条件不能保证有限极限。
\paragraph{预备定义}写 $x_n\to a$ 的完整意思是：对每个 $\eps>0$，存在整数 $K$，使所有 $n\ge K$ 满足 $|x_n-a|<\eps$。$K$ 可以依赖 $\eps$，但不能依赖随后任取的 $n$。

\Slide{27}{单调有界收敛定理的准确版本}
递增且有上界的数列收敛到全体项的上确界；递减且有下界的数列收敛到全体项的下确界。无需再额外假设递增数列下有界，因为第一项就是下界；递减情形同理。
\paragraph{证明策略}先用确界公理产生候选极限，再用确界逼近性质找到“至少一个”接近候选值的项，最后用单调性把“至少一个”升级成“此后全部”。这三个步骤是本定理的核心。
\paragraph{避免倒因为果}不能先假定极限存在再算极限。本定理正是用来证明存在的；例如递推数列先证明单调有界后，才可在递推式中取极限。

\Slide{28}{递增情形：逐句拆解证明}
设 $a=\sup\{x_n:n\in\N\}$，存在性来自非空和上有界。由上界定义，对所有 $n$ 有 $x_n\le a$。
任给 $\eps>0$，由于 $a-\eps<a$，它不可能仍是上界。因此存在某个 $K$，使 $x_K>a-\eps$。单调性告诉我们：只要 $n\ge K$，就有 $x_n\ge x_K>a-\eps$。合起来：
\[a-\eps<x_K\le x_n\le a,\qquad |x_n-a|=a-x_n<\eps.\]
这正是极限定义，故 $x_n\to a$。
\paragraph{考试易漏}确界性质只给一个接近上界的项；如果没有单调性，后面的项可能又跑远。不能把这一步省成“由确界定义显然收敛”。

\Slide{29}{递减情形与课件笔误}
\paragraph{原页勘误}证明首段的 ``supremum'' 应改为 ``infimum''；后面的 ``increasing'' 应改为 ``decreasing''。
\paragraph{完整证明}令 $a=\inf\{x_n:n\in\N\}$。非空且下有界保证存在，并且 $a\le x_n$。任给 $\eps>0$，由下确界逼近性质，存在 $K$ 使 $x_K<a+\eps$。当 $n\ge K$，递减性给 $x_n\le x_K$，因此
\[a\le x_n\le x_K<a+\eps,\qquad |x_n-a|=x_n-a<\eps.\]
故 $x_n\to a$。
\paragraph{自检}以 $x_n=1/n$ 检查：其极限是下确界 $0$，绝不是上确界 $1$。代入简单例子常能发现符号错误。

\Slide{30}{闭区间套的四种表述}
闭区间套是非空闭区间 $I_n=[d_n,u_n]$ 满足 $I_{n+1}\subseteq I_n$。于是左端点不减、右端点不增：$d_n\le d_{n+1}\le u_{n+1}\le u_n$。
结论先是交集非空，进一步精确为
\[\bigcap_{n=1}^{\infty}I_n=[D,U],\qquad D=\sup_n d_n,\quad U=\inf_n u_n.\]
若区间长度 $u_n-d_n\to0$，则 $D=U$，交集恰有一个点。
\paragraph{假设不能删除}开区间 $(0,1/n)$ 嵌套但交集为空；闭区间 $[n,\infty)$ 嵌套但不有界，交集为空；恒定区间 $[0,1]$ 说明不缩短到零就不保证唯一。

\Slide{31}{先把左端点的上确界放入每个区间}
所有 $d_n\le u_1$，因此 $D=\sup_n d_n$ 存在，并且 $d_m\le D$。还要证明 $D\le u_m$，才可说 $D\in I_m$。
\paragraph{关键交叉比较}固定任意 $m$，对任意 $n$：若 $n\le m$，则 $d_n\le d_m\le u_m$；若 $n>m$，则 $d_n\le u_n\le u_m$。所以 $u_m$ 是\emph{全部}左端点的上界，故 $D\le u_m$。
于是每个 $m$ 都有 $d_m\le D\le u_m$，即 $D\in\bigcap_m I_m$。
\paragraph{原页勘误}末尾英文 union 应为 intersection（交集）。只比较同一指标 $d_n\le u_n$ 不够，必须证明任意左右端点的交叉不等式。

\Slide{32}{再把右端点的下确界放入每个区间}
\paragraph{纠正后证明}所有 $u_n\ge d_1$，所以右端点集合\emph{下有界}，$U=\inf_n u_n$ 存在，且 $U\le u_m$。上一页的交叉比较给固定 $m$ 下，所有 $n$ 满足 $d_m\le u_n$；因此 $d_m$ 是右端点集合的下界，故 $d_m\le U$。所以 $d_m\le U\le u_m$，$U$ 属于每个区间。
\paragraph{原页勘误}bounded above 应改为 bounded below；后面的 $\sup B$ 应为 $\inf B$；$l_n$、$l_{n'}$ 应统一为左端点 $d_n,d_{n'}$；关键结论是 $d_m\le\inf B=U$，不是 $u_m\le U$。
本证明只用下确界定义，没有偷偷假设交集非空。

\Slide{33}{为什么交集恰好是一个闭区间}
要证两个集合相等，标准做法是证明两个包含关系。
\paragraph{第一方向}若 $x\in[D,U]$，每个 $n$ 都有 $d_n\le D\le x\le U\le u_n$，所以 $x\in I_n$，因此 $x$ 属于交集。
\paragraph{第二方向}若 $x\in\bigcap_n I_n$，则每个 $n$ 都有 $d_n\le x$，所以 $x$ 是全体左端点的上界，从而 $D\le x$。同理 $x\le u_n$ 对所有 $n$ 成立，使 $x$ 成为右端点集合的下界，所以 $x\le U$。故 $x\in[D,U]$。
这里的第二方向用确界定义直接完成，比课件的两次反证更短，但逻辑相同。$D\le U$ 已由左端点的上确界属于所有区间得到。

\Slide{34}{长度趋零如何产生唯一交点}
已有 $d_n\le D\le U\le u_n$，两端相减可得
\[0\le U-D\le u_n-d_n.\]
右边趋于 $0$，所以固定非负数 $U-D$ 只能为 $0$，即 $D=U=:a$。交集 $[D,U]=\{a\}$，且单调收敛定理给 $d_n\to a,u_n\to a$。
\paragraph{不引用夹逼定理也能证明}若 $U-D=c>0$，由长度趋零，某个 $n$ 有 $u_n-d_n<c$，与 $c\le u_n-d_n$ 矛盾。
\paragraph{“唯一”的证明}若 $a,b$ 都在每个区间，则 $|a-b|\le u_n-d_n$ 对所有 $n$ 成立，同样推出 $a=b$。长度条件在这一步才真正用到。

\Slide{35}{Bolzano 定理：二分法抽子列}
把全部项放进 $[D,U]$。若 $D=U$，数列常值，直接收敛。否则把区间二等分，至少一半含无穷多个\emph{项的指标}；否则两半加起来仅含有限项。保留这一半，重复二分，得到长度 $(U-D)/2^k$ 的闭区间套 $I_k$。
每个 $I_k$ 含无穷多个指标，所以可以逐次选 $n_k>n_{k-1}$ 且 $x_{n_k}\in I_k$。区间套定理给唯一交点 $a$。因为 $a,x_{n_k}$ 都在 $I_k$，
\[|x_{n_k}-a|\le |I_k|=(U-D)2^{-k}\longrightarrow0.\]
故选出的子列收敛。
\paragraph{细节}原页的严格 $<$ 一般应改为 $\le$，不影响极限。常值数列只有一个不同取值，却有无穷多个项；所以这里不能误说“无穷多个不同点”。

\Slide{36}{有界无限集合必有聚点：构造阶段}
聚点 $a$ 的定义是：每个 $\delta>0$ 都能找到 $x\in E$ 满足 $0<|x-a|<\delta$。允许 $a$ 不属于 $E$，但找出的点必须不同于 $a$。
把有界无限集合 $E$ 放入 $[-M,M]$，取 $M>0$。二分后至少一个闭半区间仍含 $E$ 的无穷多个\emph{不同点}，否则两边有限，其并也有限。反复保留无限的一半，得到 $I_n$ 嵌套，长度 $2M/2^n$，且 $I_n\cap E$ 无限。
\paragraph{与上一页的区别}数列版本保留无限指标；集合版本保留无限不同点。这一区别保证下一页能排除交点本身，真正证明“聚点”。原页 $R_0$ 应读作 $I_0$。

\Slide{37}{有界无限集合必有聚点：验证阶段}
闭区间套给唯一交点 $a$。任给 $\delta>0$，选择 $n$ 使 $|I_n|<\delta$。对每个 $x\in I_n$，因 $a\in I_n$，有 $|x-a|\le|I_n|<\delta$，因此 $I_n\subset(a-\delta,a+\delta)$。
又因为 $E\cap I_n$ 有无穷多个不同点，删去可能存在的点 $a$ 后仍至少有一个 $x$。这个 $x$ 满足 $x\in E$ 且 $0<|x-a|<\delta$。由于 $\delta$ 任意，$a$ 是聚点。
\paragraph{课件省略的关键句}只说邻域与集合相交尚不足以证明聚点，孤立点的邻域也会与集合相交。必须说“有不同于 $a$ 的点”，无限性正是用来保证这一点。

\Slide{38}{柯西列：不用预先知道极限的收敛标准}
柯西条件是
\[\forall\eps>0\ \exists K\ \forall m,n\ge K:\quad |x_n-x_m|<\eps.\]
收敛比较“项与目标”；柯西比较“后段任意两项”。因而当我们还不知道目标在哪里时，柯西条件特别有用。
\paragraph{不要只比较相邻项}$|x_{n+1}-x_n|\to0$ 比柯西弱得多。例如 $x_n=\sqrt n$ 的相邻差为 $1/(\sqrt{n+1}+\sqrt n)\to0$，但数列无界，不收敛。柯西要求任意远的两个后段项也很接近。
\paragraph{完备性的位置}收敛必为柯西在任何度量空间成立；柯西必收敛到空间内某点，是 $\R$ 的完备性，不适用于 $\Q$。

\Slide{39}{收敛推出柯西：把误差分成两半}
设 $x_n\to a$。给定最终希望得到的 $\eps>0$，按收敛定义，对精度 $\eps/2$ 选 $K$，使 $n\ge K$ 时 $|x_n-a|<\eps/2$。
若 $m,n\ge K$，利用三角不等式：
\[|x_n-x_m|=|(x_n-a)+(a-x_m)|\le|x_n-a|+|x_m-a|<\eps.\]
于是满足柯西定义。
\paragraph{为什么课件写 $2\eps$ 也能成立}因为 $\eps$ 是任意正数，可重新命名精度。但初学者考试最好从一开始取 $\eps/2$，使结尾直接得到题目要求的 $<\eps$，不留下量词解释负担。

\Slide{40}{柯西推出收敛：先有界，再锁定子列极限}
\paragraph{先证明有界}在柯西条件中先取固定精度 $1$，得 $K_0$。固定 $m=K_0$，所有 $n\ge K_0$ 满足 $|x_n|\le|x_n-x_{K_0}|+|x_{K_0}|<1+|x_{K_0}|$。再与前面有限项的绝对值取最大值，便得到全列的统一界。
\paragraph{找到候选极限}由 Bolzano 定理，有子列 $x_{n_j}\to a$。
\paragraph{证明整列趋于它}任给 $\eps>0$，柯西条件给 $K$，使 $m,n\ge K$ 时距离小于 $\eps/2$。选一个足够大的 $j$，同时保证 $n_j\ge K$ 及 $|x_{n_j}-a|<\eps/2$。以后每个 $n\ge K$ 都满足
\[|x_n-a|\le|x_n-x_{n_j}|+|x_{n_j}-a|<\eps.\]
这里先选一个共同的子列项作为桥梁，而不是对每个 $n$ 随意更换未经控制的指标。

\Slide{41}{柯西充分性的第二证法：尾部下确界}
柯西列有界的证明同上一页。定义 $y_n=\inf\{x_k:k\ge n\}$。删掉第一项后尾集缩小，下确界只能上升：$y_{n+1}\ge y_n$。原数列的统一上下界也界住 $y_n$，所以 $y_n$ 单调有界，存在极限 $a$。
\paragraph{几何解释}$y_n$ 是从第 $n$ 项起所有数的“下边界”。它不一定是实际某一项，但总有 $y_n\le x_n$。柯西条件将迫使实际项与这条下边界越来越接近。
\paragraph{原页勘误}$|nn|$ 应为 $|x_n|$。不用把下确界视作最小值；即使尾集没有最小项，定义与证明仍然有效。

\Slide{42}{尾部下确界证法：修正指标后完成}
原页混用了取下确界的哑指标与自由指标。以下使用 $j$ 表示被取遍的指标，避免混淆。
给 $\eps>0$，取 $K_1$ 使 $n\ge K_1$ 时 $|y_n-a|<\eps/2$；再取 $K_2$ 使 $n,j\ge K_2$ 时 $|x_n-x_j|<\eps/2$。
固定 $n\ge\max(K_1,K_2)$。对\emph{所有} $j\ge n$，有 $x_j>x_n-\eps/2$，所以取下确界后
\[y_n\ge x_n-\eps/2,\qquad 0\le x_n-y_n\le\eps/2.\]
从而 $|x_n-a|\le|x_n-y_n|+|y_n-a|<\eps$，故 $x_n\to a$。
\paragraph{细节}取确界后严格不等号可能变成非严格号，不能保留原页所有 $<$。由于另一误差严格小于 $\eps/2$，最终仍得到严格 $<\eps$。

\Slide{43}{第三证法：尾部上确界}
定义 $z_n=\sup\{x_k:k\ge n\}$。尾集合变小使上确界不增，故 $z_n$ 递减；柯西列有界又保证 $z_n$ 有下界。因此单调收敛定理给 $z_n\to a$。
\paragraph{对照下确界方法}下包络 $y_n$ 从下方上升，上包络 $z_n$ 从上方下降；柯西条件意味着两者之间的尾部振幅趋零。两种方法都是把不单调的数列转为单调边界。
\paragraph{原页勘误}同样应把 $|nn|$ 改为 $|x_n|$。因为 $z_n$ 是上确界，所以 $x_n\le z_n$，不能反过来。下一页用这一方向完成估计。

\Slide{44}{尾部上确界证法：正确的误差估计}
任给 $\eps>0$，取 $K$ 同时保证 $n\ge K$ 时 $|z_n-a|<\eps/2$，且 $n,j\ge K$ 时 $|x_n-x_j|<\eps/2$。
固定 $n\ge K$，对所有 $j\ge n$ 有 $x_j<x_n+\eps/2$，所以
\[x_n\le z_n=\sup_{j\ge n}x_j\le x_n+\eps/2.\]
于是 $|x_n-a|\le|x_n-z_n|+|z_n-a|<\eps$，整列收敛。
\paragraph{课件指标修正}不能把 $\sup_{j\ge K}x_j$ 写成任意的 $z_{n_1}$，两者的尾集起点不同。本页证明固定 $n$ 后统一对 $j\ge n$ 取上确界，便消除了这个问题。上下包络证法任选一种熟练掌握即可，但应认识两种思路。

\Slide{45}{开覆盖与有限子覆盖}
开集在实数轴上的意思是每个点周围都有一小段区间仍留在集合内。开覆盖 $\mathcal H$ 是一族开集，使
\[S\subseteq\bigcup_{H\in\mathcal H}H.\]
\paragraph{原页勘误}覆盖公式必须有并集符号 $\bigcup$；不能把点的集合 $S$ 直接当成开集族的子集。
有限子覆盖要求从\emph{原来给定}的覆盖中选出 $H_1,\ldots,H_m$，仍有 $S\subseteq\bigcup_{i=1}^mH_i$，不能另造一个大开集来代替。
\paragraph{Heine--Borel}实数轴中每个闭有界集合都有该性质，称为紧。$(0,1)$ 被 $(1/n,1)$（$n\ge2$）覆盖，却没有有限子覆盖，因为有限多个的并仍漏掉靠近 $0$ 的点。

\Slide{46}{有限覆盖定理：把坏区间不断二分}
先证明闭区间 $[D,U]$。反设存在一个开覆盖，无法抽出有限子覆盖。二分成两个闭半区间，至少一个仍无法有限覆盖；否则两边各取有限个，合起来仍有限，便覆盖了整个区间，矛盾。
如此递归，得到每个都“无法有限覆盖”的闭区间 $I_n$，且嵌套、长度 $(U-D)/2^n\to0$。闭区间套定理给交点 $a$。
\paragraph{思路}我们暂时不直接挑覆盖集，而是把“失败的位置”压缩到一个点。开集在该点周围会留出余地，下一页正用这个余地制造矛盾。$D=U$ 的单点区间只需选一个覆盖它的开集，已经成立。

\Slide{47}{有限覆盖定理：矛盾与一般闭集情形}
既然原族覆盖 $[D,U]$，有 $H_0\in\mathcal H$ 含 $a$。$H_0$ 开，故有 $\delta>0$ 使 $(a-\delta,a+\delta)\subseteq H_0$。取 $n$ 使 $|I_n|<\delta$。由于 $a\in I_n$，整个 $I_n$ 都在该邻域内，于是仅一个 $H_0$ 就覆盖 $I_n$，与构造矛盾。
\paragraph{补足课件未写的一步}定理针对任意闭有界 $S$，而非只有区间。取 $S\subseteq[D,U]$。因为 $S$ 闭，$\R\setminus S$ 开，把它加入原开覆盖，便覆盖整个 $[D,U]$。用已证结论取有限子覆盖，再删除其中的 $\R\setminus S$（若有），剩下的原覆盖集仍覆盖 $S$。空集情形取空覆盖即可。
\paragraph{原页记号}邻域中心须统一写交点 $a$，不能在证明中突然换成未定义的 $x$。

\Slide{48}{完备性定理的第一张复习地图}
本页复习分割、确界、单调收敛和区间套。它们的联系比孤立背诵更重要：确界帮助确定边界，单调性让后续项留在边界附近，区间套把两个单调端点结合起来。
\paragraph{逐条检查条件}分割两侧必须非空；确界存在要非空并朝相应方向有界；单调收敛要单调与有界同时满足；区间套要非空、闭、有界、嵌套，长度趋零才有唯一性。
\paragraph{原页更精确的结论}区间套交集不仅包含 $[D,U]$，而是\emph{等于} $[D,U]$。闭卷训练：不看前文写出 $d_n\le u_m$ 对任意 $m,n$ 成立的理由，这一句串起整个区间套证明。

\Slide{49}{完备性定理的第二张复习地图}
Bolzano 定理从有界数列抽出收敛子列；Weierstrass 定理为有界无限集合找聚点；柯西准则把“尾项互近”转成收敛；Heine--Borel 把任意开覆盖压缩成有限个。
\paragraph{容易混淆的结论}有界数列不一定整列收敛，如 $(-1)^n$；它只是有收敛子列。有聚点不保证集合包含该点，如 $\{1/n:n\in\N\}$。闭性才保证包含所有聚点。
\paragraph{证明工具选择}看见“有界数列”，考虑抽子列；看见“无穷多个不同点”，考虑聚点；看见“任意两项”，考虑柯西；看见一族局部邻域，考虑有限覆盖。复习时各写一个反例，检查自己是否记住假设。

\Slide{50}{闭区间上连续函数的三个定理}
连续在 $a$ 处的意思是：每给 $\eps>0$，可找到 $\delta>0$，使定义域中 $|x-a|<\delta$ 时有 $|f(x)-f(a)|<\eps$。因此 $x_n\to a$ 蕴含 $f(x_n)\to f(a)$。
\paragraph{有界性证明（子列法）}若 $f$ 在 $[D,U]$ 无界，对每个 $n$ 取 $x_n$ 使 $|f(x_n)|>n$。有界数列有子列 $x_{n_k}\to a$，且闭性给 $a\in[D,U]$。连续性给 $f(x_{n_k})\to f(a)$，但其绝对值大于 $n_k\to\infty$，矛盾。
\paragraph{有界性证明（覆盖法）}对每个 $a$，连续性给邻域 $V_a$，其内 $|f(x)|<|f(a)|+1$。有限子覆盖 $V_{a_1},\ldots,V_{a_m}$ 存在，故全局界可取 $\max_i(|f(a_i)|+1)$。
\paragraph{最大最小值证明}有界性使 $M=\sup f([D,U])$ 有限。取 $x_n$ 使 $M-1/n<f(x_n)\le M$，再抽子列趋于 $a\in[D,U]$。一方面函数值趋于 $M$，另一方面连续性使其趋于 $f(a)$，极限唯一性给 $f(a)=M$。对 $-f$ 同理得到最小值。
\paragraph{零点定理证明}不妨 $f(D)<0<f(U)$。二分区间，若中点函数值为零即结束；否则保留两端异号的半区间。若过程不停止，得到长度趋零的闭区间套，交点 $c$。左端点值始终负、右端点值始终正，连续性给 $f(c)\le0$ 且 $f(c)\ge0$，故 $f(c)=0$。端点值非零，因此 $c\in(D,U)$。

\Slide{51}{反函数连续与一致连续}
\paragraph{反函数存在}连续且严格递增的 $f:[D,U]\to\R$，由严格单调得单射，由零点定理应用于 $f-y$ 得每个 $y\in[f(D),f(U)]$ 都被取到，故有反函数 $g$。若 $y_1<y_2$ 却 $g(y_1)\ge g(y_2)$，应用递增的 $f$ 即矛盾，因此 $g$ 递增。
\paragraph{反函数连续的证明}设 $y_n\to y$。若 $g(y_n)$ 不趋于 $g(y)$，存在 $\eps_0>0$ 及子列，其每项与 $g(y)$ 距离至少 $\eps_0$。该子列有再子列趋于 $b\in[D,U]$。连续性给 $f(b)=\lim y_{n_k}=y$，单射性给 $b=g(y)$，与距离限制矛盾。故 $g$ 连续。递减情形用同样证明。
\paragraph{一致连续的量词}通常连续允许 $\delta$ 依赖点 $x$；一致连续要求一个 $\delta$ 同时控制所有点对：$|x-y|<\delta\Rightarrow |f(x)-f(y)|<\eps$。
\paragraph{一致连续证明（子列法）}若不一致连续，则存在 $\eps_0>0$，对每个 $n$ 有 $x_n,y_n\in[D,U]$ 满足 $|x_n-y_n|<1/n$ 而 $|f(x_n)-f(y_n)|\ge\eps_0$。抽取 $x_{n_k}\to a$，三角不等式给 $y_{n_k}\to a$。连续性使两函数值之差趋于零，矛盾。
\paragraph{覆盖法补充}对每个 $a$ 选 $r_a>0$，使 $|x-a|<r_a$ 时 $|f(x)-f(a)|<\eps/2$。用半径 $r_a/2$ 的邻域覆盖区间并取有限子覆盖，令 $\delta=\min_i r_{a_i}/2$。若 $|x-y|<\delta$，选含 $x$ 的小邻域，则 $x,y$ 都在对应大邻域中，用三角不等式得函数值差小于 $\eps$。

\Slide{52}{从实数到有限维向量空间}
本页重新引入向量空间，准备把一维结论推广到 $\R^N$。向量 $x=(x_1,\ldots,x_N)$ 的每个分量都是实数；加法与数乘逐坐标定义。
\paragraph{例子}$(1,2)+(3,-1)=(4,1)$，$2(1,2)=(2,4)$。在经济学中前一个坐标可表示食物数量，后一个表示衣物数量；相加表示合并两束商品。
\paragraph{结构与子集}整个 $\R^N$ 是向量空间，但非负正交锥 $\R_+^N$ 不是实向量空间，因为负数乘一个非零非负向量会跑出集合。后面讨论有界集、凸集时，不要默认它们也是向量子空间。

\Slide{53}{逐坐标验证向量公理}
实数加法交换、结合，所以对每个坐标都有 $(x+y)_i=x_i+y_i=y_i+x_i=(y+x)_i$，从而 $x+y=y+x$；其余加法公理逐坐标同理。
\paragraph{数乘分配律的示范}对每个 $i$，$[a(x+y)]_i=a(x_i+y_i)=ax_i+ay_i=[ax+ay]_i$，故 $a(x+y)=ax+ay$。等式成立的理由是所有坐标相等，不是图上看起来一致。
\paragraph{与前面公理的对应}零向量是 $(0,\ldots,0)$，逆向量是 $(-x_1,\ldots,-x_N)$。这里列出的八条与前面向量空间公理相同，作用是明确推广到多维时仍可使用熟悉的代数运算。

\Slide{54}{欧氏距离与三角不等式}
两点距离为 $\rho(x,y)=\sqrt{\sum_i(x_i-y_i)^2}$。非负性、对称性直接由平方得到；若距离为零，每个非负平方项都为零，因此每个坐标相等。
\paragraph{三角不等式的准备}Cauchy--Schwarz 不等式是 $|\sum_i u_iv_i|\le\norm u\norm v$。若 $v\ne0$，把 $t=(\sum_i u_iv_i)/\norm v^2$ 代入 $0\le\sum_i(u_i-tv_i)^2$，整理得 $(\sum_i u_iv_i)^2\le\norm u^2\norm v^2$；$v=0$ 时直接成立。
于是 $\norm{u+v}^2=\norm u^2+2\sum_i u_iv_i+\norm v^2\le(\norm u+\norm v)^2$。两边非负，开方得三角不等式；取 $u=x-y,v=y-z$ 即得距离版本。
\paragraph{距离函数连续}三角不等式交换两点后给 $|\rho(a,x)-\rho(a,y)|\le\rho(x,y)$，右边趋零就使左边趋零。这称为反三角不等式。

\Slide{55}{范数是向量的长度}
欧氏范数 $\norm x=\sqrt{\sum_i x_i^2}=\rho(x,0)$。正定性和三角不等式在上一页已证明，对称性来自 $(-x_i)^2=x_i^2$。
\paragraph{补全一般范数定义}通常还明确要求绝对齐次性 $\norm{ax}=|a|\norm x$。欧氏范数满足它，因为 $\sum_i(ax_i)^2=a^2\sum_i x_i^2$。课件只列 $\norm{-x}=\norm x$，这是 $a=-1$ 的特例，不是完整范数公理的替代。
\paragraph{多维收敛的桥梁}对每个坐标 $i$，
\[|x_i|\le\norm x\le\sqrt N\max_i|x_i|.\]
左边因单个平方项不超过总和；右边因总和至多为 $N$ 个最大平方。它说明固定有限维中，向量趋零等价于每个坐标趋零。

\Slide{56}{原页命题的反例与正确替代}
\paragraph{重要纠错}“$\R^N$ 中非空有界集合都有最大最小元素”不成立。即使 $N=1$，$(0,1)$ 非空有界，却没有最大值或最小值：对任意 $x\in(0,1)$，$(x+1)/2>x$ 且 $x/2<x$。
在 $N\ge2$ 中还必须先定义顺序。若采用逐坐标顺序，集合 $\{(1,0),(0,1)\}$ 连一个同时大于等于所有点的\emph{集合内}元素都没有，虽然它甚至有限且紧。
\paragraph{可安全使用的结论}非空有界\emph{实数集合}有上、下确界；非空紧集上的\emph{连续实值函数}达到最大与最小值。若另给连续偏好关系，极大元定理还需相应假设，不能从有界性单独推出。
本页截图保留老师原述，考试与推理中请使用带完整条件的正确命题。

\Slide{57}{多维背景下如何理解单调收敛}
本页实际陈述仍在 $\R$ 中：递增上有界趋于上确界，递减下有界趋于下确界。证明分别用确界逼近性质找到一个项，再利用单调性控制后续全部项。
\paragraph{可推广的版本}若 $x_n\in\R^N$ 在\emph{逐坐标顺序}下单调且有界，则每个坐标数列单调有界，故趋于某个 $a_i$。设 $a=(a_1,\ldots,a_N)$，由 $\norm{x_n-a}\le\sqrt N\max_i|x_{n,i}-a_i|$ 得向量收敛。
\paragraph{不要省略顺序}向量没有默认的实数式全序。需要说明“逐坐标单调”，也不能仅凭向量长度单调就推出向量收敛，例如 $(\cos n,\sin n)$ 的长度恒为 $1$。

\Slide{58}{多维闭区间套：用直径而非含糊的长度}
在 $\R^N$ 中可把区间理解为闭长方体 $I_n=\prod_{i=1}^N[d_{n,i},u_{n,i}]$。若非空、有界、嵌套，且直径 $\operatorname{diam}(I_n)=\sup\{\norm{x-y}:x,y\in I_n\}$ 趋零，则交集恰好含一点。
\paragraph{证明}每个坐标区间嵌套，其长度不超过长方体直径，故一维区间套定理给唯一 $a_i$。令 $a=(a_1,\ldots,a_N)$，则 $a\in I_n$ 对所有 $n$ 成立。若 $b$ 也是交点，$\norm{a-b}\le\operatorname{diam}(I_n)\to0$，所以 $a=b$。
\paragraph{记号说明}$\lim I_n$ 不是普通数列极限；准确说交集为 $\{a\}$，并且任取 $x_n\in I_n$ 都有 $x_n\to a$。若允许一般非空闭有界嵌套集，同样结论可由下一页的子列定理证明。

\Slide{59}{有限维 Bolzano 定理的逐坐标证明}
设 $(x_n)$ 在 $\R^N$ 有界，则每个坐标数列都有界。先对第一坐标抽收敛子列；在这个子列中对第二坐标再抽子列；继续到第 $N$ 个坐标。
\paragraph{为什么前面的收敛不会丢失}收敛数列的任何子列仍趋于原极限：原来从第 $K$ 项起误差小于 $\eps$，子列指标严格增加，最终也超过 $K$。因此最终子列的 $N$ 个坐标全部收敛。
令坐标极限为 $a_i$。给 $\eps>0$，每个坐标都能在某个指标之后误差小于 $\eps/\sqrt N$，对有限个指标取最大值，则向量误差小于 $\eps$。
\paragraph{有限维至关重要}这里只进行了有限次抽子列；不能不加条件直接搬到任意无限维空间。

\Slide{60}{有限维有界无限集合的聚点}
从有界无限集合 $E\subseteq\R^N$ 逐次选\emph{互不相同}的点 $x_1,x_2,\ldots$。有限维 Bolzano 定理给收敛子列 $x_{n_k}\to a$。
任给 $\delta>0$，充分大的 $k$ 满足 $\norm{x_{n_k}-a}<\delta$。由于所选点互异，至多一个等于 $a$，再增大一次指标即可找到不同于 $a$ 的点。因此每个去心邻域都与 $E$ 相交，$a$ 是聚点。
\paragraph{不要额外声称}$a\in E$ 未必成立；若 $E$ 闭才可以如此断言。例 $E=\{(1/n,0):n\ge1\}$ 的聚点是 $(0,0)$，它不属于 $E$。

\Slide{61}{欧氏空间的柯西准则}
若 $x_n\to a$，使用 $\norm{x_n-x_m}\le\norm{x_n-a}+\norm{x_m-a}$ 和 $\eps/2$，立即得柯西性。
\paragraph{反向逐坐标证明}若 $(x_n)$ 柯西，由 $|x_{n,i}-x_{m,i}|\le\norm{x_n-x_m}$ 可知每个坐标都是实数柯西列。实数完备性给极限 $a_i$。设 $a=(a_1,\ldots,a_N)$；给 $\eps>0$，对每个坐标选择误差阈值 $\eps/\sqrt N$，再对有限个起始指标取最大值，便得 $\norm{x_n-a}<\eps$。
这说明 $\R^N$ 是完备度量空间。逻辑上先证明实数完备性，再用有限个坐标推广，不能用尚未证明的多维结论回头证明实数结论。

\Slide{62}{有限维 Heine--Borel 定理的证明}
设闭有界 $E\subseteq\R^N$，$\mathcal H$ 是它的开覆盖。取含 $E$ 的闭正方体 $Q_0$。反设没有有限子覆盖。把正方体沿每个坐标二分，得到 $2^N$ 个小正方体；至少一个 $Q_1$ 的 $E\cap Q_1$ 不能被有限覆盖，否则把有限多份有限覆盖合并就能覆盖 $E$。
重复得到 $Q_n$ 嵌套，直径趋零，且 $E\cap Q_n$ 非空且无法有限覆盖。选 $x_n\in E\cap Q_n$；区间套定理给唯一交点 $a$，并有 $x_n\to a$。因为 $E$ 闭，$a\in E$。取 $H_0\in\mathcal H$ 含 $a$，开性给球 $B(a,r)\subseteq H_0$。当 $\operatorname{diam}(Q_n)<r$ 时，$Q_n\subseteq H_0$，矛盾。
\paragraph{为什么称为紧性}“每个开覆盖有有限子覆盖”是紧的定义。$\R^N$ 中也有逆命题：紧集有界，因为球 $B(0,m)$ 覆盖它，可抽有限个；紧集闭，因为若外点 $a$ 不在紧集，对每个 $x$ 取半径 $\norm{x-a}/3$ 的球并抽有限覆盖，有限个正距离的最小值给 $a$ 周围不碰集合的邻域。因此有限维中紧恰好等价于闭且有界。